% Generated draft: requires mathematical and source review.
\documentclass[11pt]{article}
\usepackage[margin=25mm]{geometry}
\usepackage{fontspec}
\setmainfont{Latin Modern Roman}
\usepackage{amsmath,amssymb,mathtools,unicode-math}
\setmathfont{Latin Modern Math}
\usepackage{enumitem,needspace,xurl,hyperref}
\setlength{\parindent}{0pt}
\setlength{\parskip}{5pt}
\setlength{\emergencystretch}{3em}
\newcommand{\N}{\mathbb{N}}
\newcommand{\Z}{\mathbb{Z}}
\newcommand{\Q}{\mathbb{Q}}
\newcommand{\R}{\mathbb{R}}
\newcommand{\C}{\mathbb{C}}
\newcommand{\F}{\mathbb{F}}
\newcommand{\A}{\mathbb{A}}
\newcommand{\PP}{\mathbb{P}}
\newcommand{\E}{\mathbb{E}}
\newcommand{\eps}{\varepsilon}
\newcommand{\dd}{\,\mathrm{d}}
\begin{document}
\section*{UnsolvedMath ID: 11100018}
\textbf{Source problem code:} catalogue entry
\subsection{Definitions and mathematical statement}
Fix a prime and a chromatic height $n$ in the intended Morava-theoretic setting. Let $K(n)$ be Morava $K$-theory, let $L_{K(n)}$ be Bousfield localization, and let $\mathrm{Ho}(\mathrm{Sp}_{K(n)})$ denote the homotopy category of $K(n)$-local spectra. Let $E$ be the corresponding Morava $E$-theory with coefficients $E_*=\pi_*E$. The source's notation $S$ in ``$E_*$-$S$-modules'' is not defined in the excerpt; it is customarily a Morava stabilizer group, with a possible extended-group convention, not an unqualified identification with the sphere spectrum.

\textbf{Morava $K$- and $E$-theory problem 8 (Hovey).} Understand the relationship between the $K(n)$-local category and an appropriate algebraic derived category of $E_*$-$S$-modules. This includes identifying the right algebraic category, not simply asserting an equivalence with a preselected ordinary module category.

The source notes that ordinary cooperations $E_*E=\pi_*(E\wedge E)$ form a Hopf algebroid, suggesting a comodule construction, but stresses that the relevant cooperations should instead be
\[
 \Gamma_K=\pi_*L_{K(n)}(E\wedge E).
\]
Determine the appropriate algebraic structure on these localized cooperations, its ambient category, and the resulting module/comodule and derived constructions. Completion, continuous semilinear stabilizer actions, and tensor-product conventions cannot be omitted merely by writing ordinary $E_*E$-comodules. For comparison, a customary completed homology functor is $X\mapsto\pi_*L_{K(n)}(E\wedge X)$; the excerpt does not specify an algebraic target or an equivalence theorem for it.

Hovey reports that Jens Franke claimed an equivalence and explicitly says that he did not understand the construction or even know the correct algebraic category. No prime/height hypotheses, precise Franke statement, or meaning of equivalence are supplied. A triangulated-category equivalence, a model-categorical equivalence, and an equivalence of enhanced stable categories are distinct claims. The problem as supplied asks for a relationship and a correct formulation; it does not assert all of these stronger equivalences.
\subsection{Short English statement}
Identify the correct completed algebraic module/comodule theory for Morava $E$-cooperations and explain how its derived category relates to $K(n)$-local spectra. The source reports a Franke equivalence claim without specifying its category or hypotheses.
\subsection{Sources}
\begin{itemize}
\item Local UnsolvedMath registry, permanent ID 11100018, source code catalogue entry; attributed to Mark Hovey, no proposal year supplied.
\item \emph{Hovey's Algebraic Topology Problem List}, Morava $K$- and $E$-theory, item 8; surviving mirror \url{https://www-users.cse.umn.edu/~tlawson/hovey/}. The registry records HTML-ordered-list extraction and access on 2026-07-29; the mirror has 65 top-level items on nine category pages.
\item Jens Franke is named for a claimed equivalence, but no precise paper or theorem is supplied. The summary's modern chromatic-descent interpretation also has no primary citation in the record.
\end{itemize}
\end{document}
